\documentclass[12pt, titlepage]{article}
\usepackage[colorlinks=true, linkcolor=blue, urlcolor=magenta]{hyperref}

\title{Feedback Controls Lab Weekly Report \\ \normalsize{Week 13}}
\author{Aydan Vissing, Brady Schick, Gabriel Southwell, Simon Ross}
\date{\today}

\begin{document}
\maketitle

\section*{Aydan Vissing (9 hours spent)}

After some meetings and being able to now know the dimensions of the PCB, I have been able to almost have the frame ready. Making the frame smaller took much more time than I anticipated. It will be completed by Wednesday and hopefully printed the same day. 

\section*{Brady Schick (8 hours spent)}

The time for this week was split between continuing work on the MATLAB simulation and on linearizing the mathematical model so that it can be represented with a transfer function.

\section*{Gabriel Southwell (11.5 hours spent)}

This week I got the time of flight readdressing to work!!! This process took a very long time to complete for several reasons and I am annoyed with the solution. 

Basically what I had been doing was:
\begin{itemize}
    \item Initialize sensors and store structs in arrays
    \item Use a \textit{Library} function that claims to readdress the sensor
\end{itemize}

I figured out some reasons why this did not work. When this \textit{library} function set the new address, the address it sent to the sensor was shifted by one bit. I identified this as a bug and fixed it. Unfortunately, this still did not fix it. 

Then I started to notice some other odd behavior. When I readdressed the sensors and ran an I2C scanner, they showed up on the bus with new addresses. However, when I tried to communicate with them over their new addresses, they would revert back to the default address. I spent hours with an oscilloscope trying to see what exactly was happening on the I2C bus. There were a lot of different possibilities that I checked, and I found nothing. 

I set up a meeting with Dr.\ Kohl for Monday afternoon to try and figure this out, but I ended up not needing it. This Sunday afternoon I was reading through more of the library and I found a block of code that was commented out in the init function. The code claimed to increment the I2C address of every sensor it initializes. I uncommented it and everything just worked~:/ 

Why is this not the default behavior? I do not know. Why did my method not work even though it did basically the same thing just at a different time? I do not know. 

It really does do basically the same thing that my code was supposed to. It initializes the sensor with the default address and then tells it that it has a new address and updates the struct address. 

\section*{Simon Ross (21 hours spent)}

This week was another big one for me. Basically all of it was spent tweaking the PCB design, organizing the layout, and decorating the bottom. I really enjoyed this work, so the time flew past. Dr. Kohl is ordering the final product as soon as possible, so we should have a prototype by the end of break. He said it's likely that this prototype will need some adjustment, so I think it's probable that it may be until the end of the semester before we get a working one. Hopefully we can send one home with Gabe over Christmas break.

With regards to scheduling, the sensor code is coming along well with the prototype alongside. I think I will adjust the plan to be getting a fully assembled prototype by the end of the semester. From there we can spend the next semester making it work. 

Altogether, our team is making great progress and working together well. 

\end{document}